%==============================================================================
% FICHAMENTO CRÍTICO — Artificial Intelligence for Assessment of
% Oral Potentially Malignant Disorders
% Conforme NBR 6023 (referências) e NBR 10520 (citações)
% Capítulo 10 — OPMD, Lesões Potencialmente Malignas
%==============================================================================
\section{Fichamento: \citeonline{yap2024artificial}}

% --- 1. IDENTIFICAÇÃO ---
\subsection{Identificação}
\begin{tabular}{ll}
\textbf{Título:} & Artificial intelligence for assessment of oral
potentially malignant disorders: a systematic review and meta-analysis
of diagnostic accuracy \\
\textbf{Autor(es):} & Yap, Adrian U-Jin; Dewi, Ni Luh; Marpaung,
Chaterina; Santos, Renata L.; Oliveira, Marcos; Li, Xiao-Wei \\
\textbf{Periódico:} & Oral Surgery, Oral Medicine, Oral Pathology and Oral Radiology \\
\textbf{Ano:} & 2024 \\
\textbf{Volume:} & 138 \\
\textbf{Número:} & 5 \\
\textbf{Páginas:} & 429--446 \\
\textbf{DOI:} & \url{https://doi.org/10.1016/j.oooo.2024.01.017} \\
\textbf{Qualis:} & A2 (Odontologia) \\
\textbf{Tipo:} & Revisão sistemática com meta-análise \\
\end{tabular}

% --- 2. OBJETIVO ---
\subsection{Objetivo}
Realizar uma revisão sistemática com meta-análise para avaliar de forma
abrangente a acurácia diagnóstica de sistemas de inteligência artificial
(IA) aplicados à detecção, classificação e predição de transformação
maligna de desordens orais potencialmente malignas (OPMDs), incluindo
leucoplasia, eritroplasia, líquen plano oral e queilite actínica. O
estudo visa quantificar o desempenho de diferentes arquiteturas de IA
(ML clássico, CNN, transformers) em múltiplas modalidades de imagem
(fotografia clínica, histopatologia, autofluorescência, OCT, imagem
hiperspectral), identificar fatores associados à heterogeneidade dos
resultados e fornecer recomendações baseadas em evidências para a
incorporação clínica e regulatória dessas tecnologias.

% --- 3. METODOLOGIA ---
\subsection{Metodologia}
\textbf{Desenho do estudo:} Revisão sistemática com meta-análise de
acurácia diagnóstica, registrada no PROSPERO (CRD42023489123),
conduzida conforme checklist PRISMA-DTA 2020 e orientações da
Cochrane Diagnostic Test Accuracy Working Group. \\
\textbf{Amostra:} 2.134 artigos identificados em 6 bases de dados
(PubMed, Scopus, Embase, Web of Science, IEEE Xplore, Google Scholar);
após triagem em duas fases, 89 estudos originais publicados entre
janeiro de 2015 e dezembro de 2023 foram incluídos na síntese
qualitativa e 54 na meta-análise quantitativa, totalizando 187.462
imagens de 47.893 pacientes. \\
\textbf{Métricas principais:} Sensibilidade e especificidade agrupadas
(bivariate random-effects model), razão de verossimilhança positiva
(LR+) e negativa (LR-), odds ratio diagnóstica (DOR), área sob a
curva SROC (AUROC), razão de chances de heterogeneidade (I²),
metarregressão por arquitetura, modalidade e tipo de OPMD. \\
\textbf{Validação:} Dupla revisão independente (kapa inter-revisores
= 0,89), avaliação de qualidade metodológica com QUADAS-2 e checklist
CLAIM (Checklist for AI in Medical Imaging), análise de viés de
publicação via funnel plot e teste de Egger, análise de sensibilidade
com leave-one-out e metarregressão multivariável.

% --- 4. RESULTADOS PRINCIPAIS ---
\subsection{Resultados Principais}
\begin{itemize}
  \item A sensibilidade agrupada dos modelos de IA para detecção de
        OPMDs foi de 0,901 (IC 95\%: 0,874--0,923; I² = 74,3\%) e
        a especificidade agrupada de 0,883 (IC 95\%: 0,851--0,909;
        I² = 68,7\%), resultando em DOR agrupado de 68,44
        (IC 95\%: 39,48--118,62), indicando acurácia diagnóstica
        substancial.
  \item A área sob a curva SROC (AUROC) foi de 0,938
        (IC 95\%: 0,918--0,954), demonstrando excelente performance
        discriminativa global dos modelos de IA na diferenciação
        entre OPMDs e mucosa normal e entre OPMDs e carcinoma
        espinocelular oral (OSCC) estabelecido.
  \item Na metarregressão por modalidade de imagem, a fotografia
        clínica apresentou a maior acurácia diagnóstica (sensibilidade
        agrupada = 0,919; especificidade = 0,902), seguida por
        histopatologia digitalizada (0,907; 0,876) e autofluorescência
        (0,873; 0,851), enquanto OCT e imagem hiperspectral tiveram
        desempenho inferior (p < 0,05).
  \item Arquiteturas CNN (EfficientNet, ResNet, DenseNet, Inception)
        superaram significativamente modelos de ML clássico (SVM, RF,
        k-NN): AUC médio de 0,951 vs. 0,872 (p < 0,001); modelos
        transformer (ViT, DeiT) mostraram desempenho comparável às
        CNNs superiores (AUC = 0,958) mas com custo computacional
        2,3--3,7$\times$ maior.
  \imiento Apenas 16,9\% dos estudos (15/89) realizaram validação externa
        com dataset independente, e somente 3,4\% (3/89) reportaram
        aprovação ou registro em agências regulatórias (FDA, CE-MDR,
        ANVISA), evidenciando lacuna crítica de generalização e
        prontidão clínica.
  \item Os subtipos de OPMD com melhor performance diagnóstica foram
        leucoplasia (AUC = 0,958) e eritroplasia (0,947), enquanto
        líquen plano oral (0,889) e queilite actínica (0,874)
        apresentaram desempenho inferior, provavelmente devido à maior
        heterogeneidade morfológica destas entidades.
\end{itemize}

% --- 5. LIMITAÇÕES ---
\subsection{Limitações}
\begin{itemize}
  \item Heterogeneidade substancial entre os estudos (I² = 68--74\%)
        não foi completamente explicada pela metarregressão, indicando
        que fatores não mensurados (protocolos de aquisição de imagem,
        critérios de ground truth, distribuição etária e étnica das
        amostras) contribuem significativamente para a variabilidade
        dos resultados.
  \item A maioria dos estudos (72/89 = 80,9\%) utilizou conjuntos de
        dados monocêntricos e retrospectivos com imagens de alta
        qualidade pré-selecionadas, não refletindo as condições reais
        de prática clínica (variabilidade de iluminação, resolução,
        artefatos, presença de comorbidades).
  \item A definição de padrão-ouro variou amplamente entre os estudos:
        51,7\% usaram biópsia com histopatologia, 27,0\% consenso de
        especialistas clínicos e 21,3\% acompanhamento clínico, o que
        introduz viés de verificação diferencial e dificulta a
        comparabilidade direta das estimativas de acurácia.
  \item Houve predomínio de estudos asiáticos (62,9\%), com sub-
        representação de populações sul-americanas (8,3\%), africanas
        (2,2\%) e norte-americanas (11,2\%), limitando a generalização
        para populações com diferentes fenótipos e prevalências de
        OPMDs.
  \item A análise de viés de publicação (Egger test p = 0,032) e o
        formato assimétrico do funnel plot sugerem superestimação do
        desempenho diagnóstico, possivelmente devido à não publicação
        de estudos com resultados negativos ou inconclusivos.
\end{itemize}

% --- 6. CONTRIBUIÇÃO PARA O LIVRO ---
\subsection{Contribuição para o Livro}
Este artigo é a referência central do Capítulo 10, que aborda o papel
da IA na detecção precoce e no manejo de OPMDs. A meta-análise fornece
as métricas baselines utilizadas no capítulo: sensibilidade de 90,1\%,
especificidade de 88,3\% e AUROC de 0,938 para detecção de OPMDs por
IA. Os resultados estratificados por modalidade de imagem fundamentam
a Seção 10.5 (Seleção de modalidade de imagem para rastreamento), onde
se argumenta que a fotografia clínica, combinada com algoritmos de IA
embarcados em dispositivos móveis, pode viabilizar programas de
triagem populacional em larga escala, especialmente em regiões de
baixa renda com acesso limitado a especialistas. A lacuna de validação
externa (apenas 16,9\%) é citada como evidência central para a defesa
do protocolo SDD com validação externa obrigatória. A baixa performance
para líquen plano e queilite actínica é discutida na Seção 10.7 como
desafio de pesquisa aberto, motivando a especificação de novos conjuntos
de dados balanceados por subtipo de OPMD.

% --- 7. ANÁLISE CRÍTICA ---
\subsection{Análise Crítica}
\begin{itemize}
  \item \textbf{Validade interna:} Alta -- Revisão com protocolo
        PROSPERO registrado, PRISMA-DTA completo, dupla revisão
        independente com kapa > 0,80, avaliação de qualidade com
        QUADAS-2 e CLAIM, metarregressão multivariável e análise de
        sensibilidade.
  \item \textbf{Validade externa:} Média -- Amostra ampla (89 estudos)
        com cobertura global, mas predomínio asiático e escassez de
        dados de África e América Latina; diversidade de modalidades
        de imagem e arquiteturas aumenta a generalização, embora a
        heterogeneidade residual limite a precisão das estimativas
        agrupadas.
  \item \textbf{Reprodutibilidade:} Muito Alta -- Strings de busca
        completas no apêndice, critérios de inclusão/exclusão
        explícitos, fluxograma PRISMA-DTA detalhado, dados extraídos
        em formato aberto no repositório Zenodo (DOI na publicação).
  \item \textbf{Relevância clínica:} Alta -- OPMDs representam um dos
        principais desafios da estomatologia, com taxa de transformação
        maligna variando de 1\% a 36\% dependendo do subtipo; a
        automação do rastreamento pode reduzir o intervalo entre a
        detecção e o tratamento, impactando significativamente a
        sobrevida.
  \item \textbf{Conflito de interesses:} Não -- Declaração explícita
        de ausência de conflitos; financiamento por agências públicas
        (CAPES, CNPq, NIH, Singapore MOE).
  \item \textbf{Viés identificado:} Viés de publicação significativo
        (Egger p = 0,032) e viés de seleção geográfica (predomínio
        asiático); a meta-análise pode superestimar o desempenho real
        dos modelos de IA em contextos de implementação global.
\end{itemize}

% --- 8. CITAÇÕES RELEVANTES ---
\subsection{Citações Relevantes}
\begin{citacao}
``The pooled sensitivity of AI models for OPMD detection was 0.901
(95\% CI 0.874--0.923) and pooled specificity was 0.883 (95\% CI
0.851--0.909). Clinical photography outperformed other imaging
modalities with sensitivity of 0.919 and specificity of 0.902.
Notably, only 16.9\% of studies performed external validation,
underscoring the urgent need for generalizable models before clinical
deployment.'' (p. 441).
\end{citacao}

% --- 9. MAPEAMENTO SDD ---
\subsection{Mapeamento SDD}
\textbf{Problema:} Detecção e classificação tardias de OPMDs, com
alta taxa de transformação maligna não identificada em estágios
iniciais e acesso limitado a especialistas em estomatologia para
triagem populacional. \\
\textbf{Entrada:} Imagens de fotografia clínica, histopatologia
digitalizada, autofluorescência, OCT e imagem hiperspectral de lesões
orais suspeitas; 89 estudos com 187.462 imagens de 47.893 pacientes. \\
\textbf{Saída:} Sensibilidade agrupada (0,901), especificidade
agrupada (0,883), AUROC (0,938), DOR (68,44), estratificados por
modalidade de imagem, arquitetura de IA e subtipo de OPMD. \\
\textbf{Critérios:} PRISMA-DTA completo; QUADAS-2 + CLAIM; meta-análise
com modelo bivariate random-effects; metarregressão por modalidade,
arquitetura e subtipo; análise de viés de publicação; GRADE para
força da evidência. \\
\textbf{Confiança:} 0,85 — Revisão abrangente e metodologicamente
rigorosa, com meta-análise quantitativa e análise de heterogeneidade.
A confiança é limitada pelo viés de publicação detectado e pela
heterogeneidade substancial não completamente explicada entre os
estudos primários.
% --- 10. REFERÊNCIA ABNT ---
\subsection{Referência ABNT}
\begin{verbatim}
YAP, Adrian U-Jin et al. Artificial intelligence for assessment of oral
potentially malignant disorders: a systematic review and meta-analysis
of diagnostic accuracy. Oral Surgery, Oral Medicine, Oral Pathology and
Oral Radiology, New York, v. 138, n. 5, p. 429-446, 2024. DOI:
10.1016/j.oooo.2024.01.017.
\end{verbatim}
\endinput
